%% Drafted from the mapped corpus by scripts/draft_topics.py.
% Model: gpt-5.6-terra; source characters: 20514.
\begin{konspekt}
\kreq{Формулировки}{седемте принципа --- Дирихле, биекция, събиране, изваждане, умножение, деление, включване и изключване --- \kref{2.1}.}
\kreq{Доказателство}{принцип за включване и изключване --- \kref{2.1}.}
\kreq{Дефиниции и извеждане на формулите}{четирите основни конфигурации --- с или без наредба, с или без повтаряне --- \kref{2.2}.}
\kreq{Формулировка и доказателства}{биномни коефициенти, теорема на Нютон, комбинаторни тъждества чрез двукратно броене --- \kref{2.3}.}
\kreq{Алгоритъм}{решаване на линейни рекурентни уравнения с константни коефициенти, хомогенни и нехомогенни --- \kref{2.4}.}
\kreq{Задача 1}{броене на частични функции, тотални функции, инекции и сюрекции --- \kref{2.5.1}.}
\kreq{Задача 2}{целочислени решения на $x_1+\dots+x_n=M$ с ограничения --- \kref{2.5.2}.}
\kreq{Задача 3}{решаване на зададено рекурентно уравнение до край --- \kref{2.5.3}.}
\kprereq{свойствата на бинарните релации, включително доказването и опровергаването им, принадлежат на Въпрос 1.}
\kreq{Литература}{[12], [33], [40].}
\end{konspekt}

\section{Комбинаторни принципи}\label{s:2.1}

В тази глава всички разглеждани множества са крайни, освен ако изрично не е посочено друго. Основната задача на изброителната комбинаторика е да се намери броят на обектите в дадена крайна конфигурация, като вместо непосредствено изброяване се използват структурни свойства на множествата и функциите.

\begin{keythm}[Принцип на Дирихле]
Нека $X$ и $Y$ са крайни множества и $|X|>|Y|$. Тогава не съществува инекция
\[
f:X\to Y.
\]
Еквивалентно: за всяка тотална функция $f:X\to Y$ съществуват различни $a,b\in X$, за които $f(a)=f(b)$.
\end{keythm}

\begin{proof}
Ако $f$ е инекция, всеки образ има най-много един прообраз. Непресичащите се влакна $f^{-1}(\{y\})$, $y\in Y$, покриват $X$, затова $|X|=\sum_{y\in Y}|f^{-1}(\{y\})|\le|Y|$, противоречие. Това е същият принцип, доказан във Въпрос 1.
\end{proof}

Принципът изразява, че ако повече обекти се разпределят в по-малко на брой „кутии“, поне една кутия съдържа поне два обекта.

\begin{keythm}[Принцип на биекцията]
За крайни множества $A$ и $B$ е изпълнено
\[
|A|=|B|
\]
тогава и само тогава, когато съществува биекция $f:A\to B$.
\end{keythm}

\begin{proof}
Ако $|A|=|B|=n$, избираме номерации $a:\{1,\ldots,n\}\to A$ и $b:\{1,\ldots,n\}\to B$; тогава $b\circ a^{-1}$ е биекция. Обратно, биекция $f:A\to B$ пренася всяка номерация на $A$ в номерация на $B$ със същия брой елементи. При празни множества се използва празната биекция.
\end{proof}

Принципът на биекцията е особено полезен, когато разглежданата конфигурация се съпостави еднозначно с друга, чиято мощност вече е известна.

\begin{keythm}[Принцип на събирането]
Нека
\[
A=S_1\cup S_2\cup\cdots\cup S_k,
\]
където множествата $S_1,\ldots,S_k$ са две по две непресичащи се. Тогава
\[
|A|=\sum_{i=1}^{k}|S_i|.
\]
\end{keythm}

\begin{proof}
Номерираме последователно елементите на $S_1$, после на $S_2$ и така до $S_k$. Непресичането гарантира, че няма повторения, а обединението гарантира, че не пропускаме елемент на $A$. Получената номерация има точно $\sum_i|S_i|$ позиции.
\end{proof}

Условието за непресичане е съществено: всеки елемент на $A$ трябва да бъде преброен точно веднъж.

\begin{keythm}[Принцип на изваждането]
Нека $A\subseteq U$, където $U$ е универсумът на разглежданите обекти. Тогава
\[
|A|=|U|-|U\setminus A|.
\]
\end{keythm}

\begin{proof}
Множествата $A$ и $U\setminus A$ са непресичащи се и обединението им е $U$. Принципът на събирането дава $|U|=|A|+|U\setminus A|$. Пренасяме последното събираемо от другата страна.
\end{proof}

\begin{keythm}[Принцип на умножението]
За крайни множества $A$ и $B$ е в сила
\[
|A\times B|=|A|\cdot |B|.
\]
По-общо,
\[
|X_1\times X_2\times\cdots\times X_r|
 =\prod_{i=1}^{r}|X_i|.
\]
\end{keythm}

\begin{proof}
Разбиваме $A\times B$ на непресичащите се множества $\{a\}\times B$ за $a\in A$. Всяко е в биекция с $B$ чрез $(a,b)\mapsto b$. Принципът на събирането дава $|A\times B|=\sum_{a\in A}|B|=|A||B|$. Индукция по положителния брой $r$ дава общата формула. Ако някой множител е празен, и произведението е празно.
\end{proof}

Този принцип се прилага, когато обектът се изгражда последователно чрез независим избор на компоненти.

\begin{keythm}[Принцип на делението]
Нека $A$ е непразно крайно множество и $R\subseteq A^2$ е релация на еквивалентност върху него. Ако $A$ има $k$ класа на еквивалентност и всеки клас има по $m$ елемента, то
\[
|A|=km,
\qquad\text{следователно}\qquad
m=\frac{|A|}{k}.
\]
\end{keythm}

\begin{proof}
Класовете на еквивалентност образуват разбиване. По принципа на събирането $|A|$ е сумата на размерите им, тоест $k$ еднакви събираеми $m$. Следователно $|A|=km$. Понеже разглеждаме непразно $A$, имаме $k,m>0$ и двете формули с деление са допустими.
\end{proof}

При използване за преброяване обикновено са известни $|A|$ и размерът $m$ на всеки клас; тогава броят на класовете е $|A|/m$.

\begin{keythm}[Принцип на включването и изключването]
За крайни множества $A_1,\ldots,A_n$ е изпълнено
\[
\left|\bigcup_{i=1}^{n}A_i\right|
=
\sum_{i=1}^{n}|A_i|
-\sum_{1\leq i<j\leq n}|A_i\cap A_j|
+\sum_{1\leq i<j<k\leq n}|A_i\cap A_j\cap A_k|
-\cdots
+(-1)^{n-1}|A_1\cap\cdots\cap A_n|.
\]
\end{keythm}

\begin{proof}
Доказателството е с индукция по $n$. При $n=1$ равенството е очевидно. За индукционната стъпка записваме
\[
\left|\bigcup_{i=1}^{n}A_i\right|
=
\left|\left(\bigcup_{i=1}^{n-1}A_i\right)\cup A_n\right|.
\]
За две множества е в сила
\[
|B\cup C|=|B|+|C|-|B\cap C|.
\]
Следователно
\[
\left|\bigcup_{i=1}^{n}A_i\right|
=
\left|\bigcup_{i=1}^{n-1}A_i\right|
+|A_n|
-\left|\left(\bigcup_{i=1}^{n-1}A_i\right)\cap A_n\right|.
\]
От дистрибутивния закон
\[
\left(\bigcup_{i=1}^{n-1}A_i\right)\cap A_n
=
\bigcup_{i=1}^{n-1}(A_i\cap A_n).
\]
Прилагат се индукционната хипотеза към двете обединения. След събиране на членовете се получават всички пресичания на $A_1,\ldots,A_n$ със съответните редуващи се знаци.
\end{proof}

\section{Основни комбинаторни конфигурации}\label{s:2.2}

Нека основното множество има $n$ елемента. Разглеждаме избиране на $m$ елемента, като са възможни четири случая според това дали редът е важен и дали са разрешени повторения.

\subsection{Конфигурации с наредба и повторение}\label{s:2.2.1}

\begin{defn}
Множеството от конфигурациите с наредба и повторение с дължина $m$ над множество $A$, $|A|=n$, се означава с $K_{\mathrm{нп}}(n,m)$. Неговите елементи са наредени $m$-торки
\[
(a_1,\ldots,a_m),\qquad a_i\in A.
\]
\end{defn}

\begin{thm}
\[
|K_{\mathrm{нп}}(n,m)|=n^m.
\]
\end{thm}

\begin{proof}
За всяка от $m$-те позиции има независимо по $n$ възможности. Принципът на умножението дава
\[
\underbrace{n\cdot n\cdots n}_{m\text{ пъти}}=n^m.
\]
Еквивалентно, конфигурациите са функциите от $m$-елементно множество към $n$-елементно множество.
\end{proof}

\begin{example}
Броят на думите с дължина $4$, образувани от азбука с $n$ символа, е $n^4$, ако един и същ символ може да се появява многократно.
\end{example}

\subsection{Конфигурации с наредба без повторение}\label{s:2.2.2}

\begin{defn}
Конфигурация с наредба без повторение е наредена $m$-торка от различни елементи на $n$-елементно множество. Множеството на тези конфигурации се означава с $K_{\mathrm{н}}(n,m)$.
\end{defn}

\begin{thm}
\[
|K_{\mathrm{н}}(n,m)|
=
n(n-1)\cdots(n-m+1).
\]
При $m>n$ броят е $0$.
\end{thm}

\begin{proof}
За първата позиция има $n$ възможности. След избора й остават $n-1$ възможности за втората позиция, после $n-2$ за третата и т.н. Принципът на умножението дава формулата. Ако $m>n$, инекция от $m$-елементно множество в $n$-елементно множество не съществува по принципа на Дирихле.
\end{proof}

При $m=n$ конфигурациите се наричат \emph{пермутации} и броят им е
\[
n!=n(n-1)\cdots 2\cdot 1.
\]

\subsection{Конфигурации без наредба и без повторение}\label{s:2.2.3}

\begin{defn}
Конфигурация без наредба и без повторение е $m$-елементно подмножество на $n$-елементно множество. Броят им се означава с
\[
\binom{n}{m}.
\]
\end{defn}

\begin{keythm}
За $0\leq m\leq n$ е изпълнено
\[
\binom{n}{m}
=
\frac{n(n-1)\cdots(n-m+1)}{m!}
=
\frac{n!}{m!(n-m)!}.
\]
Ако $m>n$, то
\[
\binom{n}{m}=0.
\]
\end{keythm}



\begin{proof}
Нека разгледаме множеството $K_{\mathrm{н}}(n,m)$ на наредените конфигурации без повторение. Въвеждаме релация $R$ чрез
\[
XRY
\quad\Longleftrightarrow\quad
X\text{ и }Y\text{ съдържат едни и същи елементи}.
\]
Това е релация на еквивалентност. Всеки клас на еквивалентност съдържа всички наредби на едно и също $m$-елементно множество, следователно има $m!$ елемента. Класовете са точно конфигурациите без наредба и без повторение. По принципа на делението
\[
\binom{n}{m}
=
\frac{|K_{\mathrm{н}}(n,m)|}{m!}
=
\frac{n(n-1)\cdots(n-m+1)}{m!}.
\]
\end{proof}

\subsection{Конфигурации без наредба с повторение}\label{s:2.2.4}

\begin{defn}
Конфигурация без наредба с повторение с големина $m$ над $n$-елементно основно множество е мултимножество с обща кардиналност $m$. Множеството от тези конфигурации се означава с $K_a(n,m)$.
\end{defn}

\begin{thm}
\[
|K_a(n,m)|
=
\binom{n+m-1}{m}
=
\binom{n+m-1}{n-1}.
\]
\end{thm}

\begin{proof}
Нека $x_i$ е броят на избраните обекти от $i$-тия вид. Тогава
\[
x_1+x_2+\cdots+x_n=m,
\qquad x_i\geq 0.
\]
Представяме решението чрез $m$ звездички и $n-1$ разделителя. Например броят звездички преди първия разделител е $x_1$, между първия и втория е $x_2$ и т.н. Така получаваме биекция между избора на мултимножество и избора на позициите на $m$ звездички сред общо $m+n-1$ позиции. Следователно броят е
\[
\binom{m+n-1}{m}.
\]
\end{proof}

\section{Биномни коефициенти и двукратно броене}\label{s:2.3}

\begin{prop}
За цели $0\leq m\leq n$ са в сила следните равенства; тези с избор на $1$ или $n-1$ елемента се отнасят за $n\ge1$:
\[
\binom{n}{0}=\binom{n}{n}=1,
\qquad
\binom{n}{1}=\binom{n}{n-1}=n,
\]
\[
\binom{n}{m}=\binom{n}{n-m},
\]
\[
\sum_{i=0}^{n}\binom{n}{i}=2^n.
\]
\end{prop}

\begin{proof}
Празното множество и цялото множество са единствените подмножества с размер съответно $0$ и $n$. За $n\ge1$ едноелементните подмножества са $n$, както и допълненията им. Съответствието $B\mapsto A\setminus B$ е биекция между подмножествата с размер $m$ и $n-m$. Накрая всяко от $n$-те различни елемента може независимо да бъде включено или не, затова всички подмножества са $2^n$; разбиването им по размер дава сумата.
\end{proof}

Равенството
\[
\binom{n}{m}=\binom{n}{n-m}
\]
следва от биекцията, която на всяко $m$-елементно подмножество $B$ на $A$ съпоставя допълнението $A\setminus B$, имащо $n-m$ елемента.

Равенството
\[
\sum_{i=0}^{n}\binom{n}{i}=2^n
\]
се получава, като всички подмножества на $n$-елементно множество се преброят по два начина. От една страна те са $2^n$; от друга страна ги разбиваме според кардиналността им.

\begin{keythm}[Формула на Паскал]
За $n,m\geq 1$ е в сила
\[
\binom{n}{m}
=
\binom{n-1}{m}
+
\binom{n-1}{m-1}.
\]
\end{keythm}

\begin{proof}
Нека $A$ е $n$-елементно множество и фиксираме $a\in A$. Разбиваме $m$-елементните подмножества на $A$ на две непресичащи се групи: тези, които не съдържат $a$, и тези, които съдържат $a$. В първата група има
\[
\binom{n-1}{m}
\]
подмножества. За всяко подмножество от втората група, след премахване на $a$, остава $(m-1)$-елементно подмножество на $A\setminus\{a\}$, следователно такива са
\[
\binom{n-1}{m-1}.
\]
Принципът на събирането дава търсеното равенство.
\end{proof}

\begin{keythm}[Биномна теорема на Нютон]
За $x,y\in\mathbb{R}$ и $n\in\mathbb{N}$ е изпълнено
\[
(x+y)^n
=
\sum_{k=0}^{n}\binom{n}{k}x^k y^{n-k}.
\]
\end{keythm}

\begin{proof}
Произведението $(x+y)^n$ съдържа $n$ множителя от вида $(x+y)$. При разкриване на скобите за получаване на член $x^k y^{n-k}$ трябва да изберем точно $k$ от $n$-те множителя, от които да вземем $x$; от останалите се взема $y$. Броят на тези избори е $\binom{n}{k}$. Сумирането по всички възможни $k$ дава формулата.
\end{proof}

\begin{example}
Чрез двукратно броене се получава тъждеството
\[
\binom{2n}{n}
=
\sum_{k=0}^{n}\binom{n}{k}^2.
\]
Лявата страна брои $n$-елементните подмножества на множество, разделено на две $n$-елементни части. Ако от първата част са избрани $k$ елемента, то от втората трябва да са избрани $n-k$ елемента. Броят е
\[
\binom{n}{k}\binom{n}{n-k}
=
\binom{n}{k}^2.
\]
Сумираме по всички $k$.
\end{example}

\section{Рекурентни уравнения}\label{s:2.4}

\begin{defn}
Рекурентно отношение от ред $r$ е зависимост, при която за $n\geq r$ членът $a_n$ на редица се определя чрез предходните $r$ члена:
\[
a_n=f(a_{n-1},\ldots,a_{n-r},n).
\]
Заедно с началните условия то задава редицата.
\end{defn}

\begin{example}
Отношението
\[
l(n)=
\begin{cases}
l(n-1)+n,& n\in\mathbb{N}^{*},\\
1,& n=0
\end{cases}
\]
задава числова редица. Решаването на рекурентно уравнение означава да се намери явна формула за $n$-тия член на същата редица.
\end{example}

\subsection{Хомогенни линейни рекурентни уравнения}\label{s:2.4.1}

Разглеждаме линейно рекурентно уравнение от ред $k$ с постоянни коефициенти:
\[
a_n=c_1a_{n-1}+c_2a_{n-2}+\cdots+c_ka_{n-k},
\]
където $c_k\neq 0$, а $a_0,\ldots,a_{k-1}$ са дадени начални условия.

\begin{keydefn}[Характеристично уравнение]
На горното рекурентно уравнение съответства характеристичното уравнение
\[
x^k-c_1x^{k-1}-c_2x^{k-2}-\cdots-c_k=0.
\]
\end{keydefn}

То се получава формално чрез замяна на $a_n$ с $x^n$:
\[
x^n=c_1x^{n-1}+\cdots+c_kx^{n-k},
\]
след което се дели на $x^{n-k}$.

\begin{thm}
Ако характеристичното уравнение има $k$ различни корена
\[
\alpha_1,\ldots,\alpha_k,
\]
то общото решение е
\[
a_n=A_1\alpha_1^n+\cdots+A_k\alpha_k^n,
\]
където константите $A_1,\ldots,A_k$ се определят от началните условия.
\end{thm}

\begin{proof}
Работим над $\mathbb C$ (за реална редица константите се избират така, че решението да е реално). Всяка редица $\alpha_i^n$ удовлетворява рекурентното уравнение чрез пряко заместване. Пространството на решенията е $k$-мерно, защото първите $k$ стойности се избират произволно и определят еднозначно всички следващи. Матрицата на първите $k$ стойности на тези редици е матрица на Вандермонд с детерминанта $\prod_{i<j}(\alpha_j-\alpha_i)\ne0$. Следователно редиците образуват базис и константите се определят еднозначно от началните условия.
\end{proof}

\begin{thm}
Нека различните корени на характеристичното уравнение са
\[
\beta_1,\ldots,\beta_l,
\]
като коренът $\beta_i$ има кратност $r_i$ и
\[
r_1+\cdots+r_l=k.
\]
Тогава към корена $\beta_i$ съответстват членовете
\[
A_{i,1}\beta_i^n+
A_{i,2}n\beta_i^n+
\cdots+
A_{i,r_i}n^{r_i-1}\beta_i^n.
\]
Общото решение е сумата на тези изрази за всички различни корени.
\end{thm}

\begin{proof}
Нека $E$ е операторът на изместване $(Ea)_n=a_{n+1}$ и $P(x)=\prod_i(x-\beta_i)^{r_i}$ е характеристичният полином. Уравнението е $P(E)a=0$. Понеже $c_k\ne0$, всички $\beta_i$ са ненулеви. За полином $Q$ имаме
\[(E-\beta)(\beta^nQ(n))=\beta^{n+1}\bigl(Q(n+1)-Q(n)\bigr).\]
Разликата намалява степента, затова $(E-\beta_i)^{r_i}$ анулира $\beta_i^nQ(n)$ при $\deg Q<r_i$. Тези $r_i$ редици са независими, защото ненулев полином не се анулира върху всички естествени числа. За различни $i$ пространствата са независими: при предполагаема зависимост прилагаме $\prod_{j\ne i}(E-\beta_j)^{r_j}$. Той анулира всички други слагаеми, а върху $\beta_i^nQ(n)$ е обратим, защото всяко $E-\beta_j$ има ненулев диагонален коефициент $\beta_i-\beta_j$ в полиномиалния базис. Следователно и $i$-тото слагаемо е нула. Получените общо $k$ независими решения образуват базис на $k$-мерното пространство на решенията.
\end{proof}

След намиране на общия вид се заместват дадените начални условия. Получава се система от $k$ линейни уравнения относно неизвестните константи, която се решава.

\subsection{Нехомогенни линейни рекурентни уравнения}\label{s:2.4.2}

Нехомогенно линейно рекурентно уравнение с постоянни коефициенти има вида
\[
a_n=c_1a_{n-1}+\cdots+c_ka_{n-k}
+p_1(n)b_1^n+\cdots+p_l(n)b_l^n,
\]
където $p_1(n),\ldots,p_l(n)$ са полиноми, а $b_1,\ldots,b_l$ са константи.

Първо се разглежда хомогенната част
\[
a_n=c_1a_{n-1}+\cdots+c_ka_{n-k}
\]
и се намират корените на нейното характеристично уравнение. Нехомогенната част е представена като сума от изрази от вида
\[
b^nP(n),
\]
където $P(n)$ е полином.

За частно решение на член $b^nP_d(n)$ се търси израз
\[
a_n^{(p)}=n^s b^nQ_d(n),
\]
където $Q_d$ е полином от същата степен $d$, а $s$ е кратността на $b$ като корен на характеристичния полином на хомогенната част. Ако $b$ не е корен, тогава $s=0$. Неизвестните коефициенти на $Q_d$ се намират чрез заместване и сравняване на коефициентите. За сума от такива членове се събират съответните частни решения. Общото решение е
\[
a_n=a_n^{(h)}+a_n^{(p)},
\]
а свободните константи се определят от началните условия.

\section{Примерни задачи по анотацията}\label{s:2.5}

Анотацията към този въпрос назовава три типа задачи. Методите за тях са
изложени по-горе; тук всеки тип е решен докрай, защото „методът е обяснен“ и
„задачата е решена“ са различни неща на изпита.

\subsection{Задача 1: броене на функции между крайни множества}\label{s:2.5.1}

Нека $|A|=m$ и $|B|=n$. Търсим броя на частичните функции, на тоталните
функции, на инекциите и на сюрекциите от $A$ в $B$.

\begin{prop}[Частични и тотални функции]
Броят на тоталните функции $f:A\to B$ е $n^m$, а броят на частичните функции от
$A$ в $B$ е $(n+1)^m$.
\end{prop}
\begin{proof}
Тотална функция се задава чрез независим избор на образ за всеки от $m$-те
елемента на $A$; по принципа на умножението изборите са $n^m$. Това е точно
броят на конфигурациите с наредба и повторение.

Частична функция се задава чрез избор за всеки елемент на $A$ на една от $n+1$
възможности: някой от $n$-те елемента на $B$ или „недефинирана“. Изборите отново
са независими, откъдето $(n+1)^m$. Еквивалентно: разширяваме $B$ с един нов
елемент $\bot$ и броим тоталните функции $A\to B\cup\{\bot\}$.
\end{proof}

\begin{prop}[Инекции]
Броят на инекциите $f:A\to B$ е
\[
n(n-1)\cdots(n-m+1)=\frac{n!}{(n-m)!},
\]
и е $0$ при $m>n$.
\end{prop}
\begin{proof}
Подреждаме елементите на $A$ и избираме образите последователно. Първият има $n$
възможности, вторият --- $n-1$, тъй като образът на първия е зает, и така
нататък. Това е точно броят на конфигурациите с наредба без повторение. При
$m>n$ инекция не съществува по принципа на Дирихле.
\end{proof}

\begin{prop}[Сюрекции]
Броят на сюрекциите $f:A\to B$ е
\[
\sum_{k=0}^{n}(-1)^k\binom{n}{k}(n-k)^m .
\]
\end{prop}
\begin{proof}
За $b\in B$ нека $A_b$ е множеството от тоталните функции $A\to B$, които
\emph{не} приемат стойността $b$. Сюрекциите са точно функциите извън
$\bigcup_{b\in B}A_b$. За фиксирано $k$-елементно $K\subseteq B$ пресичането
$\bigcap_{b\in K}A_b$ се състои от функциите $A\to B\setminus K$, а те са
$(n-k)^m$ на брой. Подмножествата $K$ с $k$ елемента са $\binom{n}{k}$.
Принципът за включване и изключване дава
\[
\left|\bigcup_{b\in B}A_b\right|=\sum_{k=1}^{n}(-1)^{k-1}\binom{n}{k}(n-k)^m,
\]
а изваждането от общия брой $n^m$ дава сумата от твърдението, чийто член при
$k=0$ е именно $n^m$.
\end{proof}

\begin{example}
При $m=n$ сюрекциите съвпадат с биекциите. Формулата дава
$\sum_{k=0}^{n}(-1)^k\binom{n}{k}(n-k)^n=n!$, а при $m<n$ тя дава $0$, което се
съгласува с невъзможността да се покрие по-голямо множество.
\end{example}

\subsection{Задача 2: целочислени решения с ограничения}\label{s:2.5.2}

\begin{prop}[Неотрицателни решения]
Броят на решенията на
\[
x_1+x_2+\cdots+x_n=M
\]
в неотрицателни цели числа е $\binom{M+n-1}{n-1}$.
\end{prop}
\begin{proof}
Всяко решение се кодира еднозначно с редица от $M$ еднакви звездички и $n-1$
прегради: звездичките между $(i-1)$-вата и $i$-тата преграда са стойността
$x_i$. Различните решения дават различни редици и всяка редица от $M+n-1$
позиции с избрани $n-1$ позиции за прегради дава решение. По принципа на
биекцията броят е броят на подмножествата с $n-1$ елемента, тоест
$\binom{M+n-1}{n-1}$. Това е конфигурация без наредба с повторение.
\end{proof}

\begin{prop}[Долни граници]
Нека $c_1,\ldots,c_n$ са цели числа. Броят на решенията на
$x_1+\cdots+x_n=M$ с $x_i\geq c_i$ за всяко $i$ е
\[
\binom{M-\sum_{i=1}^{n}c_i+n-1}{n-1},
\]
когато $M\geq\sum_i c_i$, и $0$ в противен случай.
\end{prop}
\begin{proof}
Полагаме $y_i=x_i-c_i$. Изображението $(x_i)\mapsto(y_i)$ е биекция между
решенията с $x_i\geq c_i$ и неотрицателните решения на
\[
y_1+\cdots+y_n=M-\sum_{i=1}^{n}c_i .
\]
Прилагаме предходното твърдение. Ако дясната страна е отрицателна, решения няма.
\end{proof}

\begin{example}[Горни граници]
Ограничение отгоре се обработва с включване и изключване, а не с изместване. За
\[
x_1+x_2+x_3=10,\qquad 0\leq x_i\leq 5,
\]
общият брой неотрицателни решения е $\binom{12}{2}=66$. Нека $B_i$ е множеството
на решенията с $x_i\geq 6$. Изместването $x_i\mapsto x_i-6$ свежда $B_i$ до
неотрицателните решения на сума $10-6=4$ с три неизвестни, тоест
$|B_i|=\binom{4+2}{2}=\binom{6}{2}=15$, а $|B_i\cap B_j|=0$, защото $6+6=12>10$.
Следователно търсеният брой е
\[
66-3\cdot 15=21 .
\]
\end{example}

\begin{example}[Смесени условия]
За $x_1+x_2+x_3+x_4=20$ с $x_1\geq 2$, $x_2\geq 3$, $x_3,x_4\geq 0$ изместваме
първите две неизвестни и получаваме $\binom{20-5+3}{3}=\binom{18}{3}=816$.
\end{example}

\subsection{Задача 3: решаване на зададено рекурентно уравнение}\label{s:2.5.3}

\begin{example}[Хомогенно уравнение с начални условия]
Да се реши
\[
a_n=5a_{n-1}-6a_{n-2},\qquad a_0=1,\quad a_1=0 .
\]
Характеристичното уравнение е $x^2-5x+6=0$ с различни корени $2$ и $3$, затова
$a_n=A\cdot 2^n+B\cdot 3^n$. Началните условия дават
\[
A+B=1,\qquad 2A+3B=0,
\]
откъдето $B=-2$ и $A=3$. Следователно
\[
a_n=3\cdot 2^n-2\cdot 3^n .
\]
Проверка: $a_0=3-2=1$, $a_1=6-6=0$ и $a_2=12-18=-6=5\cdot 0-6\cdot 1$.
\end{example}

\begin{example}[Нехомогенно уравнение с резонанс]
Да се реши
\[
a_n=3a_{n-1}-2a_{n-2}+2^n,\qquad a_0=0,\quad a_1=1 .
\]
Хомогенната част има характеристично уравнение $x^2-3x+2=0$ с корени $1$ и $2$,
тоест $a_n^{(h)}=A+B\cdot 2^n$. Основата $2$ на нехомогенния член е корен с
кратност $s=1$, а полиномът пред него е от степен $0$. Затова търсим
$a_n^{(p)}=C\,n\,2^n$. Заместваме и делим на $2^{n-2}$:
\[
4Cn=6C(n-1)-2C(n-2)+4 ,
\]
което се свежда до $0=-2C+4$, тоест $C=2$ и $a_n^{(p)}=n\cdot 2^{n+1}$. Общото
решение е
\[
a_n=A+B\cdot 2^n+n\cdot 2^{n+1}.
\]
Началните условия дават $A+B=0$ и $A+2B+4=1$, откъдето $B=-3$ и $A=3$:
\[
a_n=3-3\cdot 2^n+n\cdot 2^{n+1}.
\]
Проверка: $a_2=3-12+16=7=3\cdot 1-2\cdot 0+4$ и $a_3=3-24+48=27=3\cdot 7-2\cdot 1+8$.
\end{example}

\begin{remark}
Типична грешка е спирането при общото решение. Анотацията иска \emph{решено}
уравнение, тоест константите да са определени от началните условия и отговорът
да е проверен поне за първите няколко члена.
\end{remark}

\section{Изпитен фокус}\label{s:2.6}

\begin{itemize}
\item Формулировки на принципите на Дирихле, биекцията, събирането, изваждането, умножението, делението и включването и изключването.
\item Доказателствената схема с индукция за принципа на включването и изключването.
\item Различаване между четирите вида комбинаторни конфигурации и формулите
\[
n^m,\qquad n(n-1)\cdots(n-m+1),\qquad
\binom{n}{m},\qquad \binom{n+m-1}{m}.
\]
\item Доказване на формулата за комбинации чрез релация на еквивалентност и принципа на делението.
\item Комбинаторни доказателства на симетрията на биномните коефициенти, формулата на Паскал, равенството
\[
\sum_{i=0}^{n}\binom{n}{i}=2^n
\]
и биномната теорема на Нютон.
\item Алгоритъм за хомогенно линейно рекурентно уравнение: характеристично уравнение, различни и кратни корени, общ вид на решението и определяне на константите от началните условия.
\item Разделяне на нехомогенно рекурентно уравнение на хомогенна и нехомогенна част; за член $b^nP_d(n)$ да се търси частно решение $n^sb^nQ_d(n)$, където $s$ е кратността на $b$ като характеристичен корен.
\item Задача 1 от анотацията: броене на частични функции $(n+1)^m$, тотални функции $n^m$, инекции $n!/(n-m)!$ и сюрекции $\sum_k(-1)^k\binom{n}{k}(n-k)^m$ --- с обосновката за всяка от четирите.
\item Задача 2 от анотацията: неотрицателни целочислени решения чрез звездички и прегради; изместване $y_i=x_i-c_i$ при долни граници $x_i\geq c_i$; включване и изключване при горни граници.
\item Задача 3 от анотацията: пълно решаване на зададено хомогенно и на зададено нехомогенно уравнение --- характеристични корени, частно решение с резонансен множител $n^s$, определяне на константите от началните условия и числена проверка.
\item Свойствата на бинарните релации (рефлексивност, симетричност, транзитивност и опровергаването им) принадлежат на Въпрос 1, а не на този въпрос.
\end{itemize}
